\chapter{信息论与统计几何}
\label{chap:information-theory-and-statistical-geometry}
概率论给出了分布、条件化和换测度的语言，数理统计则把未知分布或参数作为推断对象。
不过，仅仅知道两个模型的参数相差多少，还不能判断观测能否区分它们。对
\(\mathcal{N}(0,1)\)与\(\mathcal{N}(0.1,1)\)，均值差、对数似然差、事件概率差
和搬运质量所需的距离都可以计算，但它们回答的是不同问题；若把均值重新参数化为其
立方根，参数的欧氏距离还会随坐标改变。
本章研究怎样从分布本身比较不确定性、可区分性与局部变化。编码长度引出熵，预测分布
的对数损失引出相对熵；事件、测试函数和搬运成本又给出不同的分布差异。换测度的指数
矩不等式把这些差异接到学习保证，Le Cam与Fano方法则把“分布太相近”转成任何估计器
都无法越过的误差下界。在几何一侧，Wasserstein距离从质量搬运刻画全局位置变化，Fisher度量
从KL的二阶项刻画参数模型的局部几何，变分表达把推断与分布优化连接起来。
除非另有说明，本章使用自然对数，信息量的单位为nat。若改用以\(2\)为底的对数，信息
量单位为bit，所有数值只相差常数因子。第~\ref{chap:probability-theory}章已经建立
Radon--Nikodym导数、条件分布与指数族；本章在使用这些对象时仍逐项核对绝对连续、
可积性和正则性。
\section{不确定性与熵}
\label{sec:information-entropy}
\subsection{平均编码长度为何出现对数}
设离散随机变量\(X\)取值于有限或可数字母表\(\mathcal{X}\)，质量函数为\(p(x)\)。
若使用二进制前缀码，结果\(x\)的码长记为\(\ell(x)\)。前缀码的码字不能是另一
码字的前缀，因而码长必须满足Kraft不等式
\[
  \sum_{x\in\mathcal{X}}2^{-\ell(x)}\leq1.
\]
概率较大的结果应分配较短码字，但码长不能彼此独立地任意缩短。理想实数码长
\(-\log_2p(x)\)恰好满足等号，因为\(\sum_xp(x)=1\)。
\begin{definition}[离散熵]
\label{def:information-discrete-entropy}
若\(\sum_xp(x)|\log p(x)|<\infty\)，则\(X\)的\term{熵}（entropy）定义为
\begin{equation}
  H(X)
  =
  -\sum_{x\in\mathcal{X}}p(x)\log p(x)
  =
  \mathbb{E}[-\log p(X)],
  \label{eq:information-discrete-entropy}
\end{equation}
其中约定\(0\log0=0\)。
\end{definition}
熵是按真实发生频率加权的理想描述长度。它不是某次结果携带的信息，而是结果发生前的
平均不确定性。若\(X\)几乎必然为常数，则\(H(X)=0\)；若\(X\)在\(m\)个结果上均匀
分布，则\(H(X)=\log m\)\scite{math-cover2006information}
例如，若\(X\sim\operatorname{Bernoulli}(1/4)\)，则
\[
  H(X)
  =
  -\frac14\log\frac14-\frac34\log\frac34
  \approx0.5623\ \text{nat}
  \approx0.8113\ \text{bit}.
\]
二元前缀码必须给两个结果各分配至少一位，所以其逐符号平均码长为\(1\) bit；熵给出
理想实数码长的平均值，二者的差来自整数码长约束。
\begin{theorem}[无失真前缀码的平均长度界]
\label{thm:information-source-coding-bound}
对熵有限的质量函数\(p\)和任意满足Kraft不等式的二进制码长\(\ell(x)\)，有
\[
  \mathbb{E}[\ell(X)]\geq H(X)/\log2.
\]
反过来，在有效字母表
\(\mathcal X_+=\{x:p(x)>0\}\)上取
\(\ell(x)=\lceil-\log_2p(x)\rceil\)，可构造前缀码，并满足
\[
  \mathbb{E}[\ell(X)]<H(X)/\log2+1.
\]
\end{theorem}
\begin{proof}
令\(K=\sum_x2^{-\ell(x)}\leq1\)，并在\(K>0\)时定义
\(q(x)=2^{-\ell(x)}/K\)。由\(-\log u\geq1-u\)，
\(\sum_xp(x)\log[p(x)/q(x)]\geq
1-\sum_{x:p(x)>0}q(x)\geq0\)，因此
\[
  0\leq\sum_xp(x)\log\frac{p(x)}{q(x)}
  =
  -H(X)+(\log2)\mathbb{E}[\ell(X)]+\log K.
\]
由于\(\log K\leq0\)，移项便得到下界。对取整码长，
\(2^{-\ell(x)}\leq p(x)\)，故Kraft不等式成立。为看清逆向构造，按码长从短到长排列
\(\mathcal X_+\)，依次在\([0,1)\)中放置长度为\(2^{-\ell(x)}\)且端点与该层二进
网格对齐的区间；Kraft不等式保证所有区间都能放下。用每个区间左端点的前
\(\ell(x)\)个二进制位作为码字。两个这样的二进区间若相交，较短码字必为较长码字的
前缀；构造中的区间互不相交，所以所得码是前缀码。这就是Kraft--McMillan逆向构造。
又因\(\lceil a\rceil<a+1\)，对\(p(x)>0\)逐项平均即得上界。满足\(p(x)=0\)的符号
不会由信源产生，对熵和平均码长均无贡献，因而从有效字母表中删除；若应用要求传输这些
符号，必须另外为它们保留码空间，相应冗余不由上述小于一位的界保证。
\end{proof}
这个结论只讨论无失真、逐符号可唯一解析的编码。允许块编码时，每个符号的整数取整代价
可以趋近于零；允许失真时，最优码长还取决于失真准则，不能只由\(H(X)\)决定。
\subsection{条件熵与链式法则}
对离散随机变量\(X,Y\)，给定\(Y=y\)后仍需描述\(X\)的平均长度为
\[
  H(X\mid Y=y)
  =
  -\sum_xp(x\mid y)\log p(x\mid y).
\]
再对\(Y\)平均，得到\term{条件熵}（conditional entropy）
\begin{equation}
  H(X\mid Y)
  =
  -\sum_{x,y}p(x,y)\log p(x\mid y).
  \label{eq:information-conditional-entropy}
\end{equation}
它衡量已经知道\(Y\)以后，\(X\)还剩多少不确定性。
对可数字母表，这个定义允许取值\(+\infty\)。以下推导假设
\(H(X,Y)<\infty\)，从而出现的边缘熵与条件熵均为有限值。此时
联合质量分解为\(p(x,y)=p(y)p(x\mid y)\)，所以
\[
\begin{aligned}
  H(X,Y)
  &=
  -\sum_{x,y}p(x,y)\log p(x,y)\\
  &=
  -\sum_{x,y}p(x,y)\log p(y)
  -\sum_{x,y}p(x,y)\log p(x\mid y)\\
  &=
  H(Y)+H(X\mid Y).
\end{aligned}
\]
这就是熵的链式法则。反复使用可得
\begin{equation}
  H(X_1,\ldots,X_n)
  =
  \sum_{i=1}^{n}H(X_i\mid X_1,\ldots,X_{i-1}).
  \label{eq:information-entropy-chain-rule}
\end{equation}
其中假设\(H(X_1,\ldots,X_n)<\infty\)。链式法则是同一联合分布的不同分解，
不要求变量独立。若允许扩展值，同一等式仍可由非负项求和得到，但不能把它改写成
两个无穷熵之差。
\subsection{互信息衡量观察带来的减少量}
\begin{definition}[互信息与条件互信息]
\label{def:information-mutual-information}
离散变量\(X,Y\)的\term{互信息}（mutual information, MI）定义为
\begin{equation}
  I(X;Y)
  =
  \sum_{x,y}p(x,y)
  \log\frac{p(x,y)}{p(x)p(y)}.
  \label{eq:information-mutual-information}
\end{equation}
其中\(p(x,y)=0\)的项约定为零；该对数比求和允许取值\(+\infty\)，但不拆成两个
可能无穷的熵之差。它比较实际联合质量\(p(x,y)\)与独立情形下的乘积质量
\(p(x)p(y)\)。
给定\(Z\)后的\term{条件互信息}定义为先固定\(Z=z\)计算同一对数比，再对\(Z\)平均：
\begin{equation}
  I(X;Y\mid Z)
  =
  \sum_{z:p(z)>0}p(z)
  \sum_{x,y}p(x,y\mid z)
  \log
  \frac{p(x,y\mid z)}
       {p(x\mid z)p(y\mid z)}.
  \label{eq:information-conditional-mutual-information}
\end{equation}
\end{definition}
上述定义不涉及无穷量相减。若\(H(X,Y)<\infty\)，利用熵的链式法则可进一步写成
\begin{equation}
  I(X;Y)
  =
  H(X)-H(X\mid Y)
  =
  H(Y)-H(Y\mid X).
  \label{eq:information-mi-entropy-reduction}
\end{equation}
因此它可以解释为观察一个变量后另一个变量平均减少的不确定性。条件互信息则先固定
已有信息\(Z\)，再衡量\(Y\)还能为\(X\)提供多少信息。对每个\(z\)，对数和不等式
保证内层求和非负，所以条件互信息也非负，并且可能为\(+\infty\)；在可数离散情形中，
它等于零当且仅当
\(X\perp Y\mid Z\)几乎处处成立。
\begin{theorem}[数据处理不等式]
\label{thm:information-data-processing}
若离散随机变量满足Markov关系\(X\to Y\to Z\)，即
\(X\perp Z\mid Y\)，并且相关互信息有限，则
\[
  I(X;Z)\leq I(X;Y).
\]
\end{theorem}
\begin{proof}
对同一个联合分布按两种顺序使用互信息链式法则，有
\[
\begin{aligned}
  I(X;Y,Z)
  &=I(X;Y)+I(X;Z\mid Y)=I(X;Y),\\
  I(X;Y,Z)
  &=I(X;Z)+I(X;Y\mid Z)\geq I(X;Z).
\end{aligned}
\]
第一行使用Markov关系，第二行使用条件互信息非负性。合并两式即得结论。
\end{proof}
这里的“处理”可以是确定映射，也可以是仅依赖\(Y\)的随机通道。结论说明后处理不能
凭空增加关于\(X\)的信息，但不保证任意特定任务的误差严格变差；若处理丢弃的方向与
任务无关，等号仍可能成立。一般可测空间上的版本不能依赖离散求和；下一节正式定义
KL后，将用KL的数据处理性质回收这一表达，避免有限熵中的无穷量相减。
\begin{example}[一次带噪二元观测]
\label{ex:information-binary-channel}
令\(X\)在\(\{0,1\}\)上均匀分布，观测\(Y=X\oplus N\)，其中
\(N\sim\operatorname{Bernoulli}(\varepsilon)\)与\(X\)独立，
\(0\leq\varepsilon\leq1/2\)。由于\(Y\)也均匀，
\[
  H(X)=\log2,\qquad
  H(X\mid Y)=
  -\varepsilon\log\varepsilon
  -(1-\varepsilon)\log(1-\varepsilon).
\]
记二元熵函数为
\[
  H_{\mathrm{b}}(u)
  =
  -u\log u-(1-u)\log(1-u),
  \qquad 0\leq u\leq1,
\]
其中沿用\(0\log0=0\)的约定。所以
\[
  I(X;Y)=\log2-H_{\mathrm{b}}(\varepsilon).
\]
当\(\varepsilon=0\)时，观测完整揭示\(X\)；当\(\varepsilon=1/2\)时，观测与
\(X\)独立，互信息为零。该数值描述平均可辨识信息，不保证每次观测都恢复正确。
\end{example}
\subsection{微分熵依赖坐标}
若连续随机向量\(\bm{X}\in\mathbb{R}^d\)具有密度\(p\)，并且
\(\mathbb{E}|\log p(\bm{X})|<\infty\)，定义
\term{微分熵}（differential entropy）
\[
  h(\bm{X})
  =
  -\int p(\bm{x})\log p(\bm{x})\,\mathrm{d}\bm{x}.
\]
它可以为负，也不具有离散熵“至少为零”的性质。设
\(g:\mathbb{R}^d\to\mathbb{R}^d\)是满足第~\ref{chap:probability-theory}章
密度换元定理条件的\(C^1\)微分同胚，并且
\[
  \mathbb{E}\!\left[
    \left|\log\left|\det\bm{J}_g(\bm{X})\right|\right|
  \right]
  <\infty.
\]
此时\(\bm{Y}=g(\bm{X})\)仍有密度，换元公式给出
\begin{equation}
  h(\bm{Y})
  =
  h(\bm{X})
  +
  \mathbb{E}\!\left[
    \log\left|\det\bm{J}_{g}(\bm{X})\right|
  \right].
  \label{eq:information-differential-entropy-transform}
\end{equation}
特别地，\(Y=aX\)时\(h(Y)=h(X)+\log|a|\)。坐标单位从米改成厘米就会改变微分熵，
所以不能把它单独当成坐标不变的分布差异。互信息和KL中的Jacobian矩阵项会相消，在适当
可逆变换下保持不变。
更一般地，若\(\bm{Y}=\bm{A}\bm{X}+\bm{b}\)，其中
\(\bm{A}\in\mathbb{R}^{d\times d}\)可逆，且相关积分有限，则
\[
  h(\bm{Y})=h(\bm{X})+\log|\det\bm{A}|.
\]
微分熵为负并不矛盾。例如\(X\)在区间\([0,1/2]\)上均匀分布时，密度为\(2\)，故
\(h(X)=-\log2\)。这反映密度值依赖坐标尺度，而不是出现了“负的不确定性”。
\section{预测分布与对数评分}
\label{sec:information-log-score-kl}
\subsection{交叉熵把预测质量写成平均损失}
真实分布为\(P\)，预测分布为\(Q\)。若二者在离散空间上的质量分别为\(p,q\)，观察
\(X\sim P\)后给预测者的\term{对数损失}为\(-\log q(X)\)。其平均值
\begin{equation}
  H(P,Q)
  =
  -\sum_xp(x)\log q(x)
  \label{eq:information-cross-entropy}
\end{equation}
称为\term{交叉熵}（cross-entropy）。它不是\(Q\)自身的不确定性，而是用\(Q\)
描述来自\(P\)的结果时付出的平均代价。
\begin{definition}[相对熵]
\label{def:information-kl-divergence}
同一可测空间上的概率测度\(P,Q\)若满足\(P\ll Q\)，则
\term{相对熵}（relative entropy），也称KL散度，定义为
\begin{equation}
  D_{\mathrm{KL}}(P\Vert Q)
  =
  \int\log\frac{\mathrm{d}P}{\mathrm{d}Q}\,\mathrm{d}P.
  \label{eq:information-kl-divergence}
\end{equation}
若\(P\not\ll Q\)，约定\(D_{\mathrm{KL}}(P\Vert Q)=+\infty\)。
\end{definition}
由这个定义，式~\eqref{eq:information-mutual-information}可简写为
\[
  I(X;Y)
  =
  D_{\mathrm{KL}}(P_{XY}\Vert P_X\otimes P_Y),
\]
条件互信息则是条件联合分布与条件边缘乘积分布之间KL的平均。这也把前一节的离散求和
推广到允许正则条件分布的一般可测空间。
在离散情形中，
\[
  D_{\mathrm{KL}}(P\Vert Q)
  =
  \sum_xp(x)\log\frac{p(x)}{q(x)},
\]
并且
\begin{equation}
  H(P,Q)=H(P)+D_{\mathrm{KL}}(P\Vert Q).
  \label{eq:information-cross-entropy-decomposition}
\end{equation}
真实分布固定时，最小化交叉熵等价于最小化从\(P\)到\(Q\)的KL散度。
以二值结果为例，取\(P=\operatorname{Bernoulli}(0.9)\)和
\(Q=\operatorname{Bernoulli}(0.6)\)，则
\[
\begin{aligned}
  D_{\mathrm{KL}}(P\Vert Q)&\approx0.2263,\\
  D_{\mathrm{KL}}(Q\Vert P)&\approx0.3112.
\end{aligned}
\]
两个方向使用不同分布作平均，因而数值不同。若预测分布把结果\(1\)的概率改为零，而
真实分布仍给该结果正概率，则前向交叉熵与KL都变为无穷。
\subsection{Gibbs不等式与严格适当性}
\begin{theorem}[Gibbs不等式]
\label{thm:information-gibbs-inequality}
任意概率测度\(P,Q\)满足
\[
  D_{\mathrm{KL}}(P\Vert Q)\geq0,
\]
且在散度有限时，等号成立当且仅当\(P=Q\)。
\end{theorem}
\begin{proof}
只需处理\(P\ll Q\)的情形。以\(\mu=P+Q\)为共同支配测度，密度记为\(p,q\)，
并在\(\{p>0\}\)上令\(r=q/p\)。由\(-\log u\geq1-u\)，
\[
  D_{\mathrm{KL}}(P\Vert Q)
  =
  \mathbb{E}_P[-\log r]
  \geq
  1-\mathbb{E}_P[r].
\]
若\(Q\)还在\(P\)的零集上放有质量，则
\(\mathbb{E}_P[r]\leq1\)；故右侧非负。等号要求
\(-\log r=1-r\)在\(P\)-几乎处处成立，也就是\(r=1\)，同时\(Q\)不能在
\(P\)的零集上留下额外质量。因此\(P=Q\)。
\end{proof}
式~\eqref{eq:information-cross-entropy-decomposition}与Gibbs不等式说明，对数
评分在真实分布处取得唯一最小期望损失，因而是严格适当评分规则。这是一项总体期望
性质：有限样本经验交叉熵最低的模型未必最接近真实分布，模型族限制、优化误差和样本
误差仍需分别处理。
对一般随机元素，互信息仍按
\[
  I(X;Y)=D_{\mathrm{KL}}(P_{XY}\Vert P_X\otimes P_Y)
\]
定义，条件互信息则是正则条件分布之间KL的平均；这要求所在空间允许选取相应条件分布。
Gibbs不等式立即给出非负性，并且零值分别等价于联合分布等于边缘乘积、条件联合分布
几乎处处等于条件边缘乘积。KL链式法则还给出
\[
  I(X;Y,Z)=I(X;Y)+I(X;Z\mid Y),
\]
只要相关正则条件分布存在；用扩展值形式陈述时不需要把互信息写成可能未定义的熵之差。
\subsection{方向、支持与同一高斯位置族}
KL散度通常不对称。若\(P\)在某区域有正质量而\(Q\)给出零质量，则
\(D_{\mathrm{KL}}(P\Vert Q)=+\infty\)；反向散度是否有限取决于反向支持关系。
即使两个方向都有限，数值也通常不同。因此“最小化KL”必须说明哪个分布位于前面。
KL也不满足三角不等式。令\(P,Q,R\)依次为参数\(0.1,0.2,0.3\)的Bernoulli分布，
可算得
\[
  D_{\mathrm{KL}}(P\Vert R)\approx0.1163
  >
  D_{\mathrm{KL}}(P\Vert Q)+D_{\mathrm{KL}}(Q\Vert R)
  \approx0.0624.
\]
对具有相同方差\(\sigma^2>0\)的高斯位置族
\[
  P=\mathcal{N}(\mu,\sigma^2),
  \qquad
  Q=\mathcal{N}(\nu,\sigma^2),
\]
直接代入密度比可得
\begin{equation}
  D_{\mathrm{KL}}(P\Vert Q)
  =
  \frac{(\mu-\nu)^2}{2\sigma^2}.
  \label{eq:information-gaussian-location-kl}
\end{equation}
同一计算可直接推广到向量平移。若
\[
  P=\mathcal{N}(\bm{\mu},\bm{\Sigma}),\qquad
  Q=\mathcal{N}(\bm{\nu},\bm{\Sigma}),
\]
其中\(\bm{\mu},\bm{\nu}\in\mathbb{R}^d\)，共同协方差
\(\bm{\Sigma}\in\mathbb{R}^{d\times d}\)对称正定，则
\begin{equation}
  D_{\mathrm{KL}}(P\Vert Q)
  =
  \frac12
  (\bm{\mu}-\bm{\nu})^{\top}
  \bm{\Sigma}^{-1}
  (\bm{\mu}-\bm{\nu}).
  \label{eq:information-vector-gaussian-translation-kl}
\end{equation}
这里比较的是白化后的均值位移：沿低方差方向移动相同欧氏距离更容易被观测区分。若共同
协方差只半正定，则还要检查两分布是否位于同一仿射支持上；均值差不在
\(\operatorname{range}(\bm{\Sigma})\)中时KL为无穷，不能直接把逆矩阵替换成伪逆而
省略支持条件。标量或向量公式的对称性都来自共同协方差，并不代表一般KL具有对称性。
\subsection{约束下的最大熵分布}
只给出若干矩约束时，可能有许多分布与已知信息相容。\term{最大熵原理}选择其中熵最大
者，即不额外加入约束未提供的集中结构。设参考测度为\(\mu\)，要求密度\(p\)满足
\[
  \int p\,\mathrm{d}\mu=1,
  \qquad
  \int T_jp\,\mathrm{d}\mu=\tau_j,
  \quad j=1,\ldots,k.
\]
若存在参数\(\bm{\lambda}\)使
\[
  p_{\bm{\lambda}}(x)
  =
  \exp\left\{
    \bm{\lambda}^{\top}\bm{T}(x)
    -A(\bm{\lambda})
  \right\}
\]
满足这些约束，其中
\(A(\bm{\lambda})=\log\int
e^{\bm{\lambda}^{\top}\bm{T}(x)}\,\mathrm{d}\mu(x)<\infty\)，则对任意同样可行的
密度\(p\)，
\[
\begin{aligned}
  D_{\mathrm{KL}}(p\Vert p_{\bm{\lambda}})
  &=
  \int p\log p\,\mathrm{d}\mu
  -\bm{\lambda}^{\top}\bm{\tau}
  +A(\bm{\lambda})\\
  &=
  -h(p)+h(p_{\bm{\lambda}})\geq0.
\end{aligned}
\]
所以\(p_{\bm{\lambda}}\)最大化熵，且在几乎处处意义下唯一。该结论要求可行分布和
归一化常数存在，也要求相关积分有限。矩约束不可达、自然参数落在边界或熵无上界时，
形式上的Lagrange乘子方程不保证存在最大化分布。
固定均值\(\bm{m}\)和正定协方差\(\bm{\Sigma}\)时，高斯分布
\(\mathcal{N}(\bm{m},\bm{\Sigma})\)是这一论证的代表性实例。对任意具有同样均值、
协方差且微分熵有限的密度\(p\)，高斯对数密度是关于\(\bm{x}\)的二次式，因此
\(\mathbb{E}_{p}\log q=\mathbb{E}_{q}\log q\)。于是
\[
  D_{\mathrm{KL}}(p\Vert q)
  =
  h(q)-h(p)\geq0,
\]
即
\[
  h(p)\leq
  \frac12\log\!\left((2\pi e)^d\det\bm{\Sigma}\right).
\]
该结论是在固定参考坐标和协方差约束下比较微分熵，不消除微分熵本身的坐标依赖性。
\section{分布差异的不同对象}
\label{sec:information-distribution-discrepancies}
\subsection{总变差控制事件与有界测试函数}
\begin{definition}[总变差距离]
\label{def:information-total-variation}
概率测度\(P,Q\)的\term{总变差距离}（total variation distance, TV）定义为
\begin{equation}
  d_{\mathrm{TV}}(P,Q)
  =
  \sup_{A}|P(A)-Q(A)|.
  \label{eq:information-total-variation}
\end{equation}
\end{definition}
若\(P,Q\)相对于共同测度\(\mu\)有密度\(p,q\)，令
\(A^\star=\{p\geq q\}\)，则
\begin{equation}
  d_{\mathrm{TV}}(P,Q)
  =
  \frac12\int|p-q|\,\mathrm{d}\mu
  =
  P(A^\star)-Q(A^\star).
  \label{eq:information-tv-density}
\end{equation}
证明来自\(\int(p-q)\,\mathrm{d}\mu=0\)：正部与负部积分相等，而任意事件上的差不
超过正部积分。
对任意满足\(\lVert f\rVert_\infty\leq1\)的可测函数，
\[
  |\mathbb{E}_P f-\mathbb{E}_Q f|
  \leq2d_{\mathrm{TV}}(P,Q).
\]
反过来取\(f=2\mathbb{I}_{A^\star}-1\)达到等号。因此
\begin{equation}
  2d_{\mathrm{TV}}(P,Q)
  =
  \sup_{\lVert f\rVert_\infty\leq1}
  |\mathbb{E}_P f-\mathbb{E}_Q f|.
  \label{eq:information-tv-dual}
\end{equation}
总变差同时控制全部事件概率和全部有界测试函数的期望，但对无界损失没有直接保证。
若测试函数的值域是\([0,1]\)，常数可改进为
\[
  |\mathbb{E}_P f-\mathbb{E}_Q f|
  \leq d_{\mathrm{TV}}(P,Q),
\]
因为可用层析表示\(f(x)=\int_0^1\mathbb{I}\{f(x)>t\}\,\mathrm{d}t\)，再逐层应用
事件概率差的定义。对\(P=\operatorname{Bernoulli}(1/4)\)和
\(Q=\operatorname{Bernoulli}(1/2)\)，总变差为\(1/4\)，取\(f(x)=x\)恰好达到
这个上界。
\subsection{Hellinger距离与KL控制}
取共同支配测度\(\mu\)，定义平方\term{Hellinger距离}
\begin{equation}
  d_{\mathrm{H}}^2(P,Q)
  =
  \frac12\int(\sqrt{p}-\sqrt{q})^2\,\mathrm{d}\mu
  =
  1-\int\sqrt{pq}\,\mathrm{d}\mu.
  \label{eq:information-hellinger}
\end{equation}
该值不依赖共同支配测度的选择。由
\[
  |p-q|
  =
  |\sqrt p-\sqrt q|(\sqrt p+\sqrt q)
\]
以及Cauchy--Schwarz不等式，
\[
  d_{\mathrm{H}}^2(P,Q)
  \leq
  d_{\mathrm{TV}}(P,Q)
  \leq
  \sqrt{2}\,d_{\mathrm{H}}(P,Q).
\]
左侧使用\(\sqrt p+\sqrt q\geq|\sqrt p-\sqrt q|\)，右侧使用
\(\int(\sqrt p+\sqrt q)^2\,\mathrm{d}\mu\leq4\)。
\begin{theorem}[Pinsker不等式]
\label{thm:information-pinsker}
任意概率测度\(P,Q\)满足
\begin{equation}
  d_{\mathrm{TV}}(P,Q)
  \leq
  \sqrt{\frac12D_{\mathrm{KL}}(P\Vert Q)}.
  \label{eq:information-pinsker}
\end{equation}
\end{theorem}
\begin{proof}
若KL无限，结论直接成立。否则\(P\ll Q\)。取
\(A=\{\mathrm{d}P/\mathrm{d}Q\geq1\}\)，并记
\(a=P(A)\)、\(b=Q(A)\)，则\(a-b=d_{\mathrm{TV}}(P,Q)\)。
把样本只映射成事件指标\(\mathbb{I}_A\)，对\(A\)及其补集分别应用对数和不等式，
得到KL的分组不等式
\[
  D_{\mathrm{KL}}(P\Vert Q)
  \geq
  a\log\frac{a}{b}
  +(1-a)\log\frac{1-a}{1-b}.
\]
右侧作为\(a\)的函数在\(a=b\)处函数值和一阶导数为零，二阶导数为
\(1/[a(1-a)]\geq4\)。Taylor积分余项于是给出右侧至少为
\(2(a-b)^2\)，移项即得结论。
\end{proof}
Pinsker不等式把小KL转成小事件差异，但反向一般不成立：总变差很小仍允许\(Q\)在
\(P\)有极小正质量的区域赋零，从而使KL无限。不同差异之间只能沿已证明的方向使用。
对刚才的二值分布，
\[
  d_{\mathrm{H}}^2(P,Q)\approx0.0341,\qquad
  D_{\mathrm{KL}}(P\Vert Q)\approx0.1308.
\]
相应比较给出
\(0.0341\leq0.25\leq\sqrt2\,d_{\mathrm{H}}(P,Q)\)，而Pinsker上界为
\(\sqrt{0.1308/2}\approx0.2557\)。这些数值也表明，不同差异即使可互相控制，
通常仍有不同尺度。
\subsection{测试函数类与最大均值差异}
许多任务只关心一类特定测试函数\(\mathcal{F}\)。相应的
\term{积分概率度量}（integral probability metric, IPM）为
\begin{equation}
  d_{\mathcal{F}}(P,Q)
  =
  \sup_{f\in\mathcal{F}}
  |\mathbb{E}_P f-\mathbb{E}_Q f|.
  \label{eq:information-ipm}
\end{equation}
选择所有绝对值不超过\(1\)的函数得到两倍总变差；选择所有\(1\)-Lipschitz函数将
在第~\ref{sec:information-optimal-transport}节得到一阶Wasserstein距离。
若\(\mathcal{H}\)是第~\ref{chap:functional-analysis-and-approximation}章建立的
再生核Hilbert空间，核\(k\)可测，并且
\[
  \mathbb{E}_{P}\sqrt{k(X,X)}<\infty,
  \qquad
  \mathbb{E}_{Q}\sqrt{k(Y,Y)}<\infty,
\]
则每个\(f\in\mathcal H\)在两个分布下的期望都存在。取单位球作为
\(\mathcal{F}\)，便得到\term{最大均值差异}（maximum mean discrepancy, MMD）；
它比较两个分布在该核单位球上能够产生的最大期望差。暂将相应的Riesz代表元写成核均值，
可得
\[
\begin{aligned}
  \operatorname{MMD}(P,Q)
  &=
  \sup_{\lVert f\rVert_{\mathcal{H}}\leq1}
  |\mathbb{E}_P f-\mathbb{E}_Q f|\\
  &=
  \left\lVert
    \mathbb{E}_P[k(X,\cdot)]
    -
    \mathbb{E}_Q[k(Y,\cdot)]
  \right\rVert_{\mathcal{H}}.
\end{aligned}
\]
这个等式不需要预先假定Hilbert值积分。再生性质和Cauchy--Schwarz不等式给出
\[
  |\mathbb{E}_P f(X)|
  \leq
  \lVert f\rVert_{\mathcal H}
  \mathbb{E}_P\sqrt{k(X,X)}.
\]
因此\(f\mapsto\mathbb{E}_P f(X)\)是\(\mathcal H\)上的有界线性泛函。Riesz表示
定理保证存在唯一的\(\bm{\mu}_P\in\mathcal H\)，使该泛函等于
\(\langle f,\bm{\mu}_P\rangle_{\mathcal H}\)；\(\bm{\mu}_Q\)同理。于是
\[
  \sup_{\lVert f\rVert_{\mathcal H}\leq1}
  |\langle f,\bm{\mu}_P-\bm{\mu}_Q\rangle_{\mathcal H}|
  =
  \lVert\bm{\mu}_P-\bm{\mu}_Q\rVert_{\mathcal H},
\]
达到上确界的方向在差非零时为其归一化。可测有界核自动满足上述可积条件。

核若具有特征性，MMD为零才推出\(P=Q\)；一般核可能只比较有限的一组矩或特征方向。
例如在线性核\(k(x,y)=xy\)下，核均值只记录普通均值。点质量\(\delta_0\)与在
\(\{-1,1\}\)上各取一半质量的分布均值都为零，MMD也为零，但两个分布并不相同。
\subsection{R\'enyi散度与组合边界}
对\(\alpha>1\)且\(P\ll Q\)，\term{R\'enyi散度}定义为
\begin{equation}
  D_\alpha(P\Vert Q)
  =
  \frac{1}{\alpha-1}
  \log\mathbb{E}_Q\!\left[
    \left(\frac{\mathrm{d}P}{\mathrm{d}Q}\right)^\alpha
  \right].
  \label{eq:information-renyi-divergence}
\end{equation}
它直接度量似然比的\(\alpha\)阶矩。适当条件下令\(\alpha\to1\)得到KL；阶数越大，
越强调似然比的高值区域。若\(P\not\ll Q\)，约定
\(D_\alpha(P\Vert Q)=+\infty\)；若\(P\ll Q\)但上述阶矩无限，散度也为无限。
令\(L=\mathrm{d}P/\mathrm{d}Q\)。概率空间上的\(L^p\)范数随\(p\)不减，再结合
\(\log\mathbb{E}_Q[L^\alpha]\)关于\(\alpha\)的凸性，可得
\(D_\alpha(P\Vert Q)\)随阶数不减。H\"older不等式还给出事件控制
\begin{equation}
  P(A)
  \leq
  \exp\!\left(
    \frac{\alpha-1}{\alpha}D_\alpha(P\Vert Q)
  \right)
  Q(A)^{(\alpha-1)/\alpha}.
  \label{eq:information-renyi-event-bound}
\end{equation}
这条界要求同一方向的绝对连续；若\(Q(A)=0<P(A)\)，相应R\'enyi散度必为无穷。
对独立乘积分布，密度比相乘、积分分解，故
\[
  D_\alpha(P_1\otimes P_2\Vert Q_1\otimes Q_2)
  =
  D_\alpha(P_1\Vert Q_1)+D_\alpha(P_2\Vert Q_2).
\]
这解释了为何R\'enyi散度适合记录重复随机机制的信息代价。自适应机制中，第\(t\)步
输出分布依赖历史，不能直接声称无条件独立；若对每个可能历史都有统一的条件散度上界，
才可通过条件密度比和迭代期望得到相应组合界。具体地，若每个共同历史\(h_{<t}\)下
\[
  D_\alpha\!\left(
    P_t(\cdot\mid h_{<t})\Vert Q_t(\cdot\mid h_{<t})
  \right)
  \leq\varepsilon_t,
\]
则把联合似然比写成条件似然比的乘积，从最后一步开始对其\(\alpha\)次幂作条件期望，
每次至多乘上\(e^{(\alpha-1)\varepsilon_t}\)。迭代后得到
\[
  D_\alpha(P_{1:T}\Vert Q_{1:T})
  \leq\sum_{t=1}^{T}\varepsilon_t.
\]
界对每个历史统一成立是关键；若只在某个边缘历史分布下平均，下一步自适应选择可能
偏向散度较大的历史。隐私分析使用这些结论时，还必须先固定
相邻数据定义与输出机制，不能只看到某个散度数值便推出隐私保证。
\section{换测度与信息控制}
\label{sec:information-change-measure}
\subsection{指数倾斜给出熵的变分表达}
第~\ref{sec:probability-radon-nikodym-change-measure}节说明，\(P\ll Q\)时可以用
密度比把\(P\)下期望改写为\(Q\)下期望。若只知道KL而不掌握密度比的逐点上界，指数
倾斜仍能控制换测度后的平均。
\begin{theorem}[Donsker--Varadhan变分表达]
\label{thm:information-donsker-varadhan}
设\(P\ll Q\)。对每个有界可测函数\(f\)，
\begin{equation}
  D_{\mathrm{KL}}(P\Vert Q)
  \geq
  \mathbb{E}_P[f]-\log\mathbb{E}_Q[e^f].
  \label{eq:information-dv-lower}
\end{equation}
而且在扩展实数意义下，
\begin{equation}
  D_{\mathrm{KL}}(P\Vert Q)
  =
  \sup_{f\ \text{有界可测}}
  \left\{
    \mathbb{E}_P[f]-\log\mathbb{E}_Q[e^f]
  \right\}.
  \label{eq:information-dv-variational}
\end{equation}
\end{theorem}
\begin{proof}
给定有界可测\(f\)，定义指数倾斜分布
\[
  \frac{\mathrm{d}Q_f}{\mathrm{d}Q}
  =
  \frac{e^f}{\mathbb{E}_Q[e^f]}.
\]
此时\(Q_f\)与\(Q\)等价，因而\(P\ll Q_f\)。若
\(D_{\mathrm{KL}}(P\Vert Q)<\infty\)，KL的链式相减给出
\[
\begin{aligned}
  D_{\mathrm{KL}}(P\Vert Q_f)
  &=
  \mathbb{E}_P\!\left[
    \log\frac{\mathrm{d}P}{\mathrm{d}Q}
    -f+\log\mathbb{E}_Q[e^f]
  \right]\\
  &=
  D_{\mathrm{KL}}(P\Vert Q)
  -\mathbb{E}_P[f]
  +\log\mathbb{E}_Q[e^f].
\end{aligned}
\]
左侧非负，得到式~\eqref{eq:information-dv-lower}；若左侧原KL为无穷，不等式自动
成立。

再令\(L=\mathrm{d}P/\mathrm{d}Q\)，并把\(\log0\)解释为\(-\infty\)。取有界截断
\[
  f_n=\max\{-n,\min\{\log L,n\}\}.
\]
由于\(P(L=0)=0\)，且
\(\mathbb{E}_P[(\log L)^-]
=\int_{\{L<1\}}L(-\log L)\,\mathrm{d}Q<\infty\)，截断期望满足
\(\mathbb{E}_P[f_n]\to\mathbb{E}_P[\log L]
=D_{\mathrm{KL}}(P\Vert Q)\)，右侧允许为无穷。另一方面，
\(e^{f_n}\to L\)，并且\(e^{f_n}\leq1+L\)，所以支配收敛定理给出
\(\mathbb{E}_Q[e^{f_n}]\to\mathbb{E}_Q[L]=1\)。因此对应变分目标趋于KL，证明上确界
达到所需的扩展实数值。若\(\log L\)本身有界，它可直接作为候选并取得等号；一般情形
只由上述有界截断序列逼近等号。
\end{proof}
一个有限例子可以直接核对倾斜的作用。令\(Q\)在两个点上各赋概率\(1/2\)，并取
\(f=(0,\log3)\)。归一化常数为\(\mathbb{E}_Qe^f=2\)，倾斜后
\(Q_f=(1/4,3/4)\)。直接计算可得
\[
  \mathbb{E}_{Q_f}f-D_{\mathrm{KL}}(Q_f\Vert Q)=\log2,
\]
恰好等于对数归一化常数；函数值较大的第二点从一半质量增加到四分之三。
\subsection{换测度后的期望界}
若\(g\)有界可测，把定理用于\(\lambda g\)即可得到下面的期望界。对无界函数，设
\(\lambda>0\)、\(g\in L^1(P)\)，并且
\(\mathbb{E}_Q[e^{\lambda g}]<\infty\)。令
\(g_n=\max\{-n,\min\{g,n\}\}\)，对\(\lambda g_n\)应用定理。由于
\(|g_n|\leq|g|\)且
\(e^{\lambda g_n}\leq1+e^{\lambda g}\)，分别在\(P\)和\(Q\)下使用支配收敛定理，
同样得到
\begin{equation}
  \mathbb{E}_P[g]
  \leq
  \frac{
    D_{\mathrm{KL}}(P\Vert Q)
    +\log\mathbb{E}_Q[e^{\lambda g}]
  }{\lambda}.
  \label{eq:information-change-measure-expectation}
\end{equation}
因此使用式~\eqref{eq:information-change-measure-expectation}前，必须确认
\(P\ll Q\)、\(P\)下可积性和所用参数处的指数矩有限。进一步假设
\(g\in L^1(Q)\)，并且在\(Q\)下对所有\(\lambda\in\mathbb{R}\)都有
\[
  \log\mathbb{E}_Q
  e^{\lambda(g-\mathbb{E}_Qg)}
  \leq
  \frac{\lambda^2\sigma^2}{2},
\]
则
\[
  \mathbb{E}_P[g]-\mathbb{E}_Q[g]
  \leq
  \frac{D_{\mathrm{KL}}(P\Vert Q)}{\lambda}
  +\frac{\lambda\sigma^2}{2}.
\]
对\(\lambda\)优化得到
\begin{equation}
  \mathbb{E}_P[g]-\mathbb{E}_Q[g]
  \leq
  \sqrt{2\sigma^2D_{\mathrm{KL}}(P\Vert Q)}.
  \label{eq:information-subgaussian-change-measure}
\end{equation}
这条界控制的是期望差。若需要高概率、同时对数据依赖对象成立的结论，还要把数据随机性
和候选选择纳入指数矩或联合控制。
\subsection{信息稳定性与PAC--Bayes接口}
设样本\(S\)经过随机学习算法产生输出\(W\)。取
\[
  P=P_{S,W},
  \qquad
  Q=P_S\otimes P_W,
\]
则换测度代价恰为\(I(S;W)\)。如果某个中心化泛化差
\(g(S,w)\)对每个固定\(w\)在\(P_S\)下具有统一次高斯指数矩，式
\eqref{eq:information-subgaussian-change-measure}便把平均泛化差控制为互信息的
平方根。这里的信息量衡量算法输出对样本的依赖，不等同于参数个数。
PAC--Bayes使用相同代数，但比较对象不同。设\(P_0\)是见数据前固定的参数先验，
\(Q_S\)是数据依赖的随机预测分布。对满足上一小节可积性和指数矩条件的参数函数
\(f(\theta)\)，
\[
  \mathbb{E}_{\theta\sim Q_S}[f(\theta)]
  \leq
  D_{\mathrm{KL}}(Q_S\Vert P_0)
  +
  \log\mathbb{E}_{\theta\sim P_0}[e^{f(\theta)}].
\]
再对样本建立指数矩和高概率控制，才得到具体PAC--Bayes风险界。一个确定性的换测度
不等式本身不是泛化定理；先验是否独立于评价样本、损失是否有界或次高斯、后验如何由
数据选择，都必须在具体结果中说明。

下面给出一个可以直接回用的基础版本。设
\(S=(Z_1,\ldots,Z_n)\)独立同分布，损失
\(\ell(\theta,Z)\in[0,1]\)，并定义
\[
  R(\theta)=\mathbb{E}[\ell(\theta,Z)],
  \qquad
  \widehat R_S(\theta)=\frac1n\sum_{i=1}^{n}\ell(\theta,Z_i).
\]
令\(P_0\)与\(S\)独立，并在观察样本前固定。对任意预先固定的
\(\lambda>0\)和\(0<\delta<1\)，以至少\(1-\delta\)的概率，同时对所有
\(Q\ll P_0\)成立
\begin{equation}
  \mathbb{E}_{Q}R(\theta)
  \leq
  \mathbb{E}_{Q}\widehat R_S(\theta)
  +
  \frac{
    D_{\mathrm{KL}}(Q\Vert P_0)+\log(1/\delta)
  }{\lambda}
  +
  \frac{\lambda}{8n}.
  \label{eq:information-pac-bayes-basic}
\end{equation}
事实上，对每个固定\(\theta\)，Hoeffding引理给出
\[
  \mathbb{E}_S
  e^{\lambda(R(\theta)-\widehat R_S(\theta))}
  \leq e^{\lambda^2/(8n)}.
\]
先对\(P_0\)平均，再用Markov不等式可知，在同一个概率至少为\(1-\delta\)的事件上，
\[
  \mathbb{E}_{P_0}
  e^{\lambda(R-\widehat R_S)}
  \leq
  \frac{e^{\lambda^2/(8n)}}{\delta}.
\]
在这个事件内应用Donsker--Varadhan不等式，便得到式
\eqref{eq:information-pac-bayes-basic}。事件只涉及先验积分，所以结论能同时覆盖
事件发生后按数据选择的全部\(Q\)。若再用数据选择\(\lambda\)或\(P_0\)，则需离散
并集界、样本拆分或新的统一指数矩，不能直接沿用这条同时性。

同一损失设定也给出信息稳定性的具体常数。若随机算法输出\(W\)，则对固定\(w\)，
\(R(w)-\widehat R_S(w)\)在\(P_S\)下满足方差代理\(1/(4n)\)的次高斯界。将联合分布
\(P_{S,W}\)与边缘乘积比较，式~\eqref{eq:information-subgaussian-change-measure}
给出
\[
  \left|
    \mathbb{E}\bigl[R(W)-\widehat R_S(W)\bigr]
  \right|
  \leq
  \sqrt{\frac{I(S;W)}{2n}}.
\]
这是平均的有符号泛化差界，不是对每个样本或每个算法输出都成立的高概率保证。
\subsection{轨迹KL的链式分解}
考虑随机序列\(Z_{0:T}=(Z_0,\ldots,Z_T)\)。设\(P,Q\)都按历史分解为
\[
  p(z_{0:T})
  =
  p_0(z_0)\prod_{t=1}^{T}p_t(z_t\mid z_{<t}),
\]
以及对应的\(q\)，并假设每个\(P\)条件核在相关历史上绝对连续于\(Q\)条件核。取
对数似然比后求\(P\)下期望，塔式性质给出
\begin{equation}
\begin{aligned}
  D_{\mathrm{KL}}(P_{0:T}\Vert Q_{0:T})
  &=
  D_{\mathrm{KL}}(P_0\Vert Q_0)\\
  &\quad+
  \sum_{t=1}^{T}
  \mathbb{E}_{Z_{<t}\sim P}
  \left[
    D_{\mathrm{KL}}\!\left(
      P_t(\cdot\mid Z_{<t})
      \Vert
      Q_t(\cdot\mid Z_{<t})
    \right)
  \right].
  \label{eq:information-trajectory-kl-chain}
\end{aligned}
\end{equation}
这不是独立样本的乘积公式。每一步的分布可以依赖完整历史，但条件KL必须沿\(P\)实际
诱导的历史分布平均。
在序贯决策中，轨迹分布由环境转移与算法策略共同决定。若比较两个环境而保持同一算法，
算法在同一历史上的条件动作核相同，可以在似然比中相消；但两个环境产生的历史频率仍
不同，期望必须沿相应轨迹分布计算。若策略也改变，动作核的KL项不能省略。
例如，在共同环境中比较两条历史依赖策略时，环境条件核逐项相消，剩余项是沿
\(P\)轨迹平均的条件动作KL之和；在共同策略下比较两个环境时，动作核相消，剩余项
则按该策略实际访问到的状态动作对累计环境条件KL。访问次数因此是随机量，其期望由
参考环境与算法共同决定，不能预先替换成均匀的独立样本数。

式~\eqref{eq:information-trajectory-kl-chain}直接适用于固定有限时域。若在停时
\(\tau\leq T\)停止，可以在停止后补上共同的吸收符号，把问题化为长度\(T\)的轨迹；
此时停止后的条件KL为零。无界停时则还需保证几乎处处终止，并对累计对数似然比建立
可积控制或一致可积性，才能把有限截断的等式取极限。
\section{信息论下界与难以区分的模型}
\label{sec:information-lower-bounds}
\subsection{二元检验的最小错误由总变差决定}
设观测\(Z\)来自\(P_0\)或\(P_1\)，两个假设先验各为\(1/2\)。检验
\(\psi(Z)\in\{0,1\}\)把某个事件\(A=\{\psi=1\}\)判为假设\(1\)。两类错误之和为
\[
  P_0(A)+P_1(A^{\mathsf c})
  =
  1-\{P_1(A)-P_0(A)\}.
\]
对\(A\)优化并使用总变差定义，得到
\begin{equation}
  \inf_\psi
  \{P_0(\psi=1)+P_1(\psi=0)\}
  =
  1-d_{\mathrm{TV}}(P_0,P_1).
  \label{eq:information-binary-testing-tv}
\end{equation}
这里的左侧是两类条件错误概率之和；在两个假设先验各为\(1/2\)时，最小总体错误概率为
\[
  \frac12\left\{1-d_{\mathrm{TV}}(P_0,P_1)\right\}.
\]
所以分布的总变差越小，任何检验都越难同时降低两类错误。
\subsection{Le Cam两点方法}
\begin{theorem}[Le Cam两点下界]
\label{thm:information-le-cam}
设两个模型\(P_0,P_1\)对应目标\(\theta_0,\theta_1\)，并在度量\(d\)下满足
\(d(\theta_0,\theta_1)\geq2s\)。对任意估计量\(\widehat{\theta}\)，
\begin{equation}
  \max_{j\in\{0,1\}}
  P_j\!\left(
    d(\widehat{\theta},\theta_j)\geq s
  \right)
  \geq
  \frac{1-d_{\mathrm{TV}}(P_0,P_1)}{2}.
  \label{eq:information-le-cam}
\end{equation}
\end{theorem}
\begin{proof}
由估计量构造检验
\[
  \psi=
  \mathbb{I}\{
    d(\widehat{\theta},\theta_1)
    <
    d(\widehat{\theta},\theta_0)
  \}.
\]
若真实模型为\(0\)但\(\psi=1\)，三角不等式说明
\(d(\widehat{\theta},\theta_0)\geq s\)；否则估计值同时离\(\theta_0\)小于\(s\)
且离\(\theta_1\)更近，会与两点距离至少\(2s\)矛盾。模型\(1\)同理。因此两个估计
坏事件的概率和至少为任意二元检验的错误和。使用式
\eqref{eq:information-binary-testing-tv}，再以最大值不小于平均值，即得结论。
\end{proof}
若观察\(n\)个独立同分布样本，则
\[
  D_{\mathrm{KL}}(P_0^{\otimes n}\Vert P_1^{\otimes n})
  =
  nD_{\mathrm{KL}}(P_0\Vert P_1).
\]
结合Pinsker不等式，可以用单样本KL控制Le Cam下界。例如对
\(\mathcal{N}(-\delta,\sigma^2)\)与
\(\mathcal{N}(\delta,\sigma^2)\)，目标均值相距\(2\delta\)，而\(n\)样本KL为
\(2n\delta^2/\sigma^2\)。取\(\delta\)与\(\sigma/\sqrt n\)同阶，可使两模型仍
不能可靠区分，从而得到均值估计误差至少为\(n^{-1/2}\)量级。这个构造给出下界量级，
不声称任意常数都最优。
例如取\(\delta=\sigma/(2\sqrt n)\)，则\(n\)样本KL等于\(1/2\)，Pinsker不等式
给出总变差至多\(1/2\)。式~\eqref{eq:information-le-cam}于是说明，对任何均值
估计量，至少在一个模型下，以不小于\(1/4\)的概率产生不小于
\(\sigma/(2\sqrt n)\)的误差。这里常数来自所选两点和Pinsker放缩，并非最优常数。
\subsection{Fano不等式处理多个候选}
两点方法不能反映候选数随问题规模增长的困难。令随机索引
\(V\)在\(\{1,\ldots,M\}\)上均匀，给定\(V=v\)后观测
\(Z\sim P_v\)。任意解码器\(\widehat V(Z)\)的错误概率记为
\(p_{\mathrm e}=\mathbb{P}(\widehat V\neq V)\)
\scite{math-cover2006information}
\begin{theorem}[Fano不等式]
\label{thm:information-fano}
当\(M\geq2\)时，
\begin{equation}
  p_{\mathrm e}
  \geq
  1-\frac{I(V;Z)+\log2}{\log M}.
  \label{eq:information-fano}
\end{equation}
\end{theorem}
\begin{proof}
令\(E=\mathbb{I}\{\widehat V\neq V\}\)。由于\(\widehat V\)是\(Z\)的函数，
\[
\begin{aligned}
  H(V\mid Z)
  &\leq H(V,E\mid Z)\\
  &=H(E\mid Z)+H(V\mid E,Z)\\
  &\leq\log2+p_{\mathrm e}\log(M-1)\\
  &\leq\log2+p_{\mathrm e}\log M.
\end{aligned}
\]
另一方面，\(V\)均匀使
\(H(V\mid Z)=H(V)-I(V;Z)=\log M-I(V;Z)\)。合并两式并移项即得结论。
\end{proof}
为了把Fano不等式变成估计下界，可在参数空间中选取一个\(2s\)-分离的有限打包
\(\{\theta_1,\ldots,\theta_M\}\)。若估计误差小于\(s\)，最近邻解码就能恢复索引；
所以解码错误蕴含估计误差至少为\(s\)。随后还需控制\(I(V;Z)\)，常用上界为
\[
  I(V;Z)
  =
  \frac1M\sum_{v=1}^{M}
  D_{\mathrm{KL}}(P_v\Vert\overline P)
  \leq
  \frac1{M^2}\sum_{v,u}
  D_{\mathrm{KL}}(P_v\Vert P_u),
\]
其中\(\overline P=M^{-1}\sum_uP_u\)。打包既要让目标相隔足够远，又要让观测分布的
KL足够小；只完成其中一项不能得到有效下界。
一个一维数值构造已经能展示两项约束如何配合。取\(M=4\)，令候选均值为
\(\{-3a,-a,a,3a\}\)，每个候选产生\(n\)个方差为\(\sigma^2\)的独立高斯观测。
这些均值两两至少相隔\(2a\)，而有序候选对的平均平方距离为\(10a^2\)，所以
\[
  I(V;Z)
  \leq
  \frac{5na^2}{\sigma^2}.
\]
取\(a=\sigma/(4\sqrt n)\)，右侧为\(5/16\)。代入Fano不等式可得
\(p_{\mathrm e}\gtrsim0.274\)，从而任意均值估计器都在某个候选下以非零常数概率
产生至少\(a\)的误差。高维问题通常用更多分离方向使\(\log M\)随维数增长，但仍须
同时核对参数间损失与观测分布间KL。
\subsection{自适应访问下的轨迹下界}
在强化学习或主动实验中，第\(t\)次观测位置由过去数据决定。此时不能把样本写成固定
设计下的独立乘积。应对算法与环境共同诱导的轨迹分布使用式
\eqref{eq:information-trajectory-kl-chain}。若两个环境只在某些状态动作对上不同，
轨迹KL通常具有
\[
  D_{\mathrm{KL}}(\mathbb{P}_0^\mathcal{A}
  \Vert\mathbb{P}_1^\mathcal{A})
  =
  \sum_{s,a}
  \mathbb{E}_0[N_T(s,a)]
  D_{\mathrm{KL}}\!\left(
    P_0(\cdot\mid s,a)\Vert P_1(\cdot\mid s,a)
  \right)
\]
这样的访问次数分解。这里\(N_T(s,a)\)本身由算法和环境决定。
要推出样本复杂度或遗憾下界，还必须让两个环境的最优行动、奖励或价值确实相差一定量。
仅构造轨迹分布接近的环境不够，因为它们可能要求完全相同的决策；仅构造价值不同的环境
也不够，因为算法可能用一次观测就区分它们。信息下界的核心是同时满足“目标分离”和
“观测接近”。
例如有两个可选择来源\(A,B\)。环境\(0\)中二者回报均值分别为
\(1/2+\varepsilon\)和\(1/2\)，环境\(1\)只把\(B\)的均值改为
\(1/2+2\varepsilon\)，其中\(0<\varepsilon<1/4\)。两个环境的最优来源相反，
但全部区分信息只来自访问\(B\)时的Bernoulli观测。因而任意自适应算法的轨迹KL为
\[
  \mathbb{E}_0[N_T(B)]
  D_{\mathrm{KL}}\!\left(
    \operatorname{Bernoulli}(1/2)
    \Vert
    \operatorname{Bernoulli}(1/2+2\varepsilon)
  \right).
\]
若算法在两个环境中都以小错误概率选出最优来源，二元检验下界要求这个轨迹KL不能太小，
从而迫使\(\mathbb{E}_0[N_T(B)]\)达到相应量级。有限时域可直接使用该式；随机无界
停止时刻还需沿用前节的终止与可积条件，不能无条件把\(T\)换成\(\tau\)。
\section{最优传输与分布距离}
\label{sec:information-optimal-transport}
\subsection{耦合描述质量怎样搬运}
总变差关心同一事件上的质量差，KL关心同一点上的密度比。若两个分布形状相同但位置略有
平移，它们的支持甚至可能不重叠，此时KL可为无限，而“把质量移动多远”仍然很小。
\term{最优传输}（optimal transport, OT）正是寻找这种最小搬运代价
\scite{math-villani2009transport}
给定空间\(\mathcal{X},\mathcal{Y}\)上的概率测度\(P,Q\)，
\term{耦合}（coupling）是边缘分别为\(P,Q\)的联合分布。全体耦合记为
\(\Pi(P,Q)\)。对非负成本\(c(x,y)\)，Kantorovich最优传输代价为
\begin{equation}
  \mathcal{T}_c(P,Q)
  =
  \inf_{\gamma\in\Pi(P,Q)}
  \int c(x,y)\,\mathrm{d}\gamma(x,y).
  \label{eq:information-kantorovich-transport}
\end{equation}
耦合决定来源质量与目标质量如何配对；优化在所有满足边缘约束的配对中选择成本最低者。
在一般空间上，最小值是否达到需要额外条件。若\(\mathcal X,\mathcal Y\)是Polish
空间，也就是完备可分度量空间，成本\(c\)非负且下半连续，并且至少存在一个有限成本
耦合，则存在有限成本的最优耦合。下面给出证明所用的拓扑概率接口。概率测度称为
\term{紧的}（tight），是指对任意\(\varepsilon>0\)都能找到紧集，使其外部质量小于
\(\varepsilon\)；一族测度若可共用这个紧集，就称为\term{一致紧的}（uniformly
tight）。Polish空间上的每个概率测度都紧，因此可分别选择紧集
\(K_{\mathcal X},K_{\mathcal Y}\)，使任意\(\gamma\in\Pi(P,Q)\)都满足
\[
  \gamma((K_{\mathcal X}\times K_{\mathcal Y})^{\mathsf c})
  \leq
  P(K_{\mathcal X}^{\mathsf c})+Q(K_{\mathcal Y}^{\mathsf c})
  <\varepsilon.
\]
所以全部耦合，特别是一列近最优耦合，构成一致紧族。Prokhorov定理把一致紧性转成相对
弱序列紧性，因而可取子列\(\gamma_k\Rightarrow\gamma\)；这里弱收敛表示对每个有界
连续函数\(h\)，都有\(\int h\,\mathrm{d}\gamma_k\to\int h\,\mathrm{d}\gamma\)。
把\(h\)取为只依赖一个坐标的有界连续函数，可知弱极限仍保持边缘\(P,Q\)，所以
\(\gamma\in\Pi(P,Q)\)。最后，Portmanteau定理对非负下半连续成本给出
\[
  \int c\,\mathrm{d}\gamma
  \leq
  \liminf_{k\to\infty}\int c\,\mathrm{d}\gamma_k.
\]
这里变化的是测度，不能直接把固定测度下的Fatou引理当作这一步的证明。
在有限集合上，若来源质量为\(p_i\)、目标质量为\(q_j\)，成本为\(c_{i,j}\)，问题是
\[
  \min_{\gamma_{i,j}\geq0}
  \sum_{i,j}c_{i,j}\gamma_{i,j},
  \quad
  \sum_j\gamma_{i,j}=p_i,\quad
  \sum_i\gamma_{i,j}=q_j.
\]
这是第~\ref{chap:optimization}章线性规划的直接实例。其对偶为
\begin{equation}
  \max_{\bm{u},\bm{v}}
  \left\{
    \sum_i u_ip_i+\sum_jv_jq_j:
    u_i+v_j\leq c_{i,j}
  \right\}.
  \label{eq:information-discrete-ot-dual}
\end{equation}
这个约束可以从原问题完整消去搬运变量得到。分别用\(u_i\)和\(v_j\)乘两组边缘
等式，并保留\(\gamma_{i,j}\geq0\)作为取下确界时的定义域，Lagrangian为
\begin{align}
  \mathcal{L}(\bm{\gamma},\bm{u},\bm{v})
  &=
  \sum_{i,j}c_{i,j}\gamma_{i,j}
  +\sum_i u_i\left(p_i-\sum_j\gamma_{i,j}\right)
  +\sum_j v_j\left(q_j-\sum_i\gamma_{i,j}\right)
  \nonumber\\
  &=
  \sum_i u_ip_i+\sum_jv_jq_j
  +\sum_{i,j}(c_{i,j}-u_i-v_j)\gamma_{i,j}.
  \label{eq:information-discrete-ot-lagrangian}
\end{align}
现在逐个消去\(\gamma_{i,j}\)。若某个系数
\(c_{i,j}-u_i-v_j<0\)，令相应\(\gamma_{i,j}\to+\infty\)便使
\(\inf_{\bm{\gamma}\geq0}\mathcal{L}=-\infty\)；若全部系数非负，下确界在
\(\bm{\gamma}=\bm0\)处达到，并等于
\(\sum_i u_ip_i+\sum_jv_jq_j\)。因此对偶函数为
\[
  g(\bm{u},\bm{v})
  =
  \begin{cases}
    \sum_i u_ip_i+\sum_jv_jq_j,
      &u_i+v_j\leq c_{i,j}\text{ 对所有 }i,j,\\
    -\infty,&\text{否则}.
  \end{cases}
\]
最大化这个对偶函数，正好得到式~\eqref{eq:information-discrete-ot-dual}。
约束\(u_i+v_j\leq c_{i,j}\)不是额外猜出的条件，而是保证消去非负搬运变量后
下确界有限的必要且充分条件。弱对偶也可由该约束逐项乘以
\(\gamma_{i,j}\)并求和得到；有限维可行且最优值有限时，线性规划强对偶保证两者相等。
有限情形的存在性更直接：约束集合非空，因为独立耦合
\(\gamma_{i,j}=p_iq_j\)总是可行；它又是闭且有界的有限维多面体，所以线性目标能在
其上取到最小值。耦合不必来自确定映射，一个源点的质量可以分送到多个目标点。

以位置\(0,2\)上的两个分布为例，取\(\bm p=(3/4,1/4)\)、
\(\bm q=(1/4,3/4)\)，成本\(c_{i,j}=|x_i-y_j|\)。可行搬运表
\[
  \bm{\gamma}
  =
  \begin{pmatrix}
    1/4 & 1/2\\
    0   & 1/4
  \end{pmatrix}
\]
的成本为\(1\)。对偶变量\(\bm u=(0,-2)\)、\(\bm v=(0,2)\)满足全部
\(u_i+v_j\leq c_{i,j}\)，且对偶值同样为
\(\bm u^{\top}\bm p+\bm v^{\top}\bm q=1\)。弱对偶说明任何可行搬运的成本至少为
\(1\)，所以上述表和势函数共同构成可核对的最优性证书。
\subsection{Wasserstein距离与矩条件}
本小节假设\((\mathcal{X},d)\)是Polish度量空间。固定\(x_0\in\mathcal X\)，并将
具有有限\(r\)阶矩的概率测度记为
\[
  \mathcal P_r(\mathcal X)
  =
  \left\{
    P:\int d(x,x_0)^r\,\mathrm{d}P(x)<\infty
  \right\};
\]
这个集合不依赖基点\(x_0\)的选择。对\(r\geq1\)和
\(P,Q\in\mathcal P_r(\mathcal X)\)，定义
\begin{equation}
  W_r(P,Q)
  =
  \left[
    \inf_{\gamma\in\Pi(P,Q)}
    \int d(x,y)^r\,\mathrm{d}\gamma(x,y)
  \right]^{1/r}.
  \label{eq:information-wasserstein-distance}
\end{equation}
在\(\mathcal P_r(\mathcal X)\)上，\(W_r\)构成距离。三角不等式的关键是使用
Polish空间上的粘合引理，把
\(P,Q\)的耦合与\(Q,R\)的耦合沿共同边缘\(Q\)粘合成三元联合分布，再对
\(d(X,Z)\leq d(X,Y)+d(Y,Z)\)使用Minkowski不等式。若相应矩无限，距离可以为
无限。
这些距离性质使用了底层\(d\)的非负、对称、点分离和三角不等式。一般自定义成本
\(c(x,y)\)得到的\(\mathcal T_c\)仍是搬运代价，但不自动对称，也不自动满足三角
不等式，因而不能只因使用了耦合优化就称为距离。
对\(r=1\)，在上述Polish空间和有限一阶矩条件下，
Kantorovich--Rubinstein对偶给出
\begin{equation}
  W_1(P,Q)
  =
  \sup_{\operatorname{Lip}(f)\leq1}
  |\mathbb{E}_P f-\mathbb{E}_Q f|.
  \label{eq:information-kr-duality}
\end{equation}
因此若任务损失\(f\)关于输入是\(L\)-Lipschitz，就有
\[
  |\mathbb{E}_P f-\mathbb{E}_Q f|
  \leq
  L W_1(P,Q).
\]
成本度量必须与任务敏感方向相符；在无关坐标上代价很大，可能使Wasserstein距离很大却
不影响任务，反之亦然。
\subsection{同一高斯位置族下的三种比较}
仍取
\[
  P=\mathcal{N}(\mu,\sigma^2),
  \qquad
  Q=\mathcal{N}(\nu,\sigma^2).
\]
令\(X\sim P\)，并取\(Y=X+\nu-\mu\)，得到一个耦合，其位移恒为
\(|\mu-\nu|\)。另一方面，任意耦合都满足
\[
  \mathbb{E}|X-Y|^2
  \geq
  |\mathbb{E}X-\mathbb{E}Y|^2
  =
  |\mu-\nu|^2.
\]
所以
\begin{equation}
  W_2(P,Q)=|\mu-\nu|.
  \label{eq:information-gaussian-location-wasserstein}
\end{equation}
与式~\eqref{eq:information-gaussian-location-kl}相比，Wasserstein距离不除以
噪声尺度，KL则衡量平移相对于观测噪声有多容易辨认。
例如当\(|\mu-\nu|=0.2\)、\(\sigma=1\)时，
\begin{align*}
  W_2(P,Q)&=0.2,\\
  D_{\mathrm{KL}}(P\Vert Q)&=0.02,\\
  d_{\mathrm{TV}}(P,Q)&=2\Phi(0.1)-1\approx0.0797,
\end{align*}
其中\(\Phi\)为标准正态分布函数。若均值差保持不变而把\(\sigma\)减半，
\(W_2\)不变，KL增至\(0.08\)，总变差也随可辨识性增大。
若把两个点质量\(\delta_0,\delta_\varepsilon\)比较，则
\(W_r(\delta_0,\delta_\varepsilon)=|\varepsilon|\)，但两个方向的KL在
\(\varepsilon\neq0\)时都为无限，因为支持不兼容。总变差则恒为\(1\)，不随
\(\varepsilon\to0\)而减小。这说明三种差异的拓扑和任务含义不同，不能只按数值大小
互换。
\subsection{分布演化与比较边界}
一族随时间变化的分布\((P_t)\)可以在Wasserstein空间中形成路径，路径长度反映质量随
时间移动的累计代价。这为平均场模型、粒子系统和连续生成过程提供几何语言。但边缘分布
之间距离较小，不自动说明产生这些边缘的轨迹机制接近；路径依赖、条件转移和时间耦合
仍需单独建模。
MMD与一阶Wasserstein距离都可写成测试函数类差异，却选择不同函数类。MMD的计算和
统计性质依赖核，\(W_1\)的含义依赖底层成本；高阶Wasserstein距离仍由搬运代价定义，
并不具有这里的IPM表达。高维经验Wasserstein距离还可能收敛缓慢。选择哪一个不是纯粹
的数学偏好，应从任务损失、支持结构、可用矩条件和估计代价共同判断。
\section{Fisher几何与自然梯度}
\label{sec:information-fisher-geometry}
\subsection{KL的局部二阶项定义度量}
设\(\{P_{\bm{\theta}}:\bm{\theta}\in\Theta\subseteq\mathbb{R}^d\}\)是由共同测度
支配的参数模型，密度为\(p_{\bm{\theta}}(x)\)。固定内点\(\bm{\theta}\)，并假设存在
其开邻域\(U\subseteq\Theta\)与共同可测支持\(S\)，使所有
\(\bm{\eta}\in U\)都满足\(p_{\bm{\eta}}>0\)于\(S\)、在\(S\)外为零。再假设对几乎
处处的\(x\in S\)，映射
\(\bm{\eta}\mapsto p_{\bm{\eta}}(x)\)二阶连续可微，且其一、二阶导数分别由
\(\mu\)-可积函数在\(U\)上一致控制。定义得分向量
\[
  \bm{s}_{\bm{\theta}}(x)
  =
  \nabla_{\bm{\theta}}\log p_{\bm{\theta}}(x).
\]
这些条件允许把一、二阶微分移入积分。由
\(\int p_{\bm{\theta}}\,\mathrm{d}\mu=1\)，
\begin{equation}
  \mathbb{E}_{\bm{\theta}}[
    \bm{s}_{\bm{\theta}}(X)
  ]
  =
  \int\nabla_{\bm{\theta}}p_{\bm{\theta}}(x)\,\mathrm{d}\mu(x)
  =
  \bm{0}.
  \label{eq:information-score-zero}
\end{equation}
为使下面两种Fisher表示都成为有限矩阵，进一步假设
\[
  \mathbb{E}_{\bm{\theta}}
  \lVert\bm{s}_{\bm{\theta}}(X)\rVert_2^2
  <\infty,
  \qquad
  \mathbb{E}_{\bm{\theta}}\!\left[
    \left\lVert
      \nabla_{\bm{\theta}}^2\log p_{\bm{\theta}}(X)
    \right\rVert_{\mathrm{op}}
  \right]
  <\infty.
\]
\begin{definition}[Fisher信息矩阵]
\label{def:information-fisher-matrix}
模型在\(\bm{\theta}\)处的\term{Fisher信息矩阵}定义为
\begin{equation}
  \bm{F}(\bm{\theta})
  =
  \mathbb{E}_{\bm{\theta}}\!\left[
    \bm{s}_{\bm{\theta}}(X)
    \bm{s}_{\bm{\theta}}(X)^{\top}
  \right].
  \label{eq:information-fisher-matrix}
\end{equation}
\end{definition}
关键代数不能只把式~\eqref{eq:information-score-zero}中的得分当作唯一依赖参数的
对象。由乘积法则，
\[
  \nabla_{\bm{\theta}}^2p_{\bm{\theta}}
  =
  p_{\bm{\theta}}\!\left(
    \nabla_{\bm{\theta}}^2\log p_{\bm{\theta}}
    +
    \bm{s}_{\bm{\theta}}\bm{s}_{\bm{\theta}}^{\top}
  \right).
\]
在上述可积性和交换微分条件下，对两边积分并使用
\(\int\nabla_{\bm{\theta}}^2p_{\bm{\theta}}\,\mathrm{d}\mu=\bm0\)，得到信息恒等式
\[
  \bm{F}(\bm{\theta})
  =
  -\mathbb{E}_{\bm{\theta}}\!\left[
    \nabla_{\bm{\theta}}^2
    \log p_{\bm{\theta}}(X)
  \right].
\]
前一个表达保证半正定，后一个表达把Fisher信息与期望对数似然曲率相连。支持随参数
改变时，边界项可能破坏信息恒等式。
\begin{theorem}[KL的局部二阶展开]
\label{thm:information-kl-local-fisher}
除上述共同支持与密度导数条件外，假设\(U\)可取为凸集，
\(\bm{F}(\bm{\theta})\)的所有元素有限，并且对
\(P_{\bm{\theta}}\)-几乎处处的\(x\)，
\(\nabla_{\bm{\eta}}^2\log p_{\bm{\eta}}(x)\)在
\(\bm{\eta}=\bm{\theta}\)处连续。进一步要求存在
\(M\in L^1(P_{\bm{\theta}})\)，使
\[
  \sup_{\bm{\eta}\in U}
  \left\lVert
    \nabla_{\bm{\eta}}^2\log p_{\bm{\eta}}(x)
  \right\rVert_{\mathrm{op}}
  \leq M(x)
\]
几乎处处成立。则当\(\bm{\delta}\to\bm{0}\)且
\(\bm{\theta}+\bm{\delta}\in U\)时，
\begin{equation}
  D_{\mathrm{KL}}\!\left(
    P_{\bm{\theta}}
    \Vert
    P_{\bm{\theta}+\bm{\delta}}
  \right)
  =
  \frac12
  \bm{\delta}^{\top}
  \bm{F}(\bm{\theta})
  \bm{\delta}
  +
  o(\lVert\bm{\delta}\rVert_2^2).
  \label{eq:information-kl-fisher-expansion}
\end{equation}
\end{theorem}
\begin{proof}
对固定\(x\)，沿线段
\(\bm{\theta}+t\bm{\delta}\)使用带积分余项的Taylor公式：
\[
\begin{aligned}
  \log p_{\bm{\theta}+\bm{\delta}}(x)
  &=
  \log p_{\bm{\theta}}(x)
  +
  \bm{\delta}^{\top}\bm{s}_{\bm{\theta}}(x)\\
  &\quad+
  \int_0^1(1-t)\,
  \bm{\delta}^{\top}
  \nabla_{\bm{\eta}}^2
    \log p_{\bm{\eta}}(x)
    \big|_{\bm{\eta}=\bm{\theta}+t\bm{\delta}}
  \bm{\delta}\,\mathrm{d}t.
\end{aligned}
\]
把两边移入
\(\mathbb{E}_{\bm{\theta}}[
\log p_{\bm{\theta}}-\log p_{\bm{\theta}+\bm{\delta}}]\)。
一阶项由得分均值为零而消失。Hessian的逐点连续性和\(M\)的可积控制允许使用支配收敛
定理，得到
\[
  \int_0^1(1-t)\,
  \mathbb{E}_{\bm{\theta}}\!\left[
    \nabla_{\bm{\eta}}^2\log p_{\bm{\eta}}(X)
    \big|_{\bm{\eta}=\bm{\theta}+t\bm{\delta}}
  \right]\mathrm{d}t
  =
  \frac12
  \mathbb{E}_{\bm{\theta}}\!\left[
    \nabla_{\bm{\theta}}^2\log p_{\bm{\theta}}(X)
  \right]+o(1).
\]
再使用信息恒等式
\(\mathbb{E}_{\bm{\theta}}[
\nabla_{\bm{\theta}}^2\log p_{\bm{\theta}}(X)]
=-\bm{F}(\bm{\theta})\)，并在两侧乘入
\(\bm{\delta}\)，便得到二阶项及
\(o(\lVert\bm{\delta}\rVert_2^2)\)余项。
\end{proof}
要让这个二次型成为切空间上的局部内积，还需假设得分映射满秩：对任意
\(\bm{v}\neq\bm{0}\)，都有
\(\mathbb{E}_{\bm{\theta}}[
(\bm{v}^{\top}\bm{s}_{\bm{\theta}}(X))^2]>0\)。
这等价于\(\bm{F}(\bm{\theta})\)正定，也可表述为参数到分布的映射在该点是正则浸入。
局部可识别本身并不足够。例如一维高斯位置族
\(P_\theta=\mathcal{N}(\theta^3,\sigma^2)\)的参数映射一一对应，但
\[
  F(\theta)=\frac{9\theta^4}{\sigma^2},
\]
所以它在\(\theta=0\)处局部可识别却没有正定Fisher信息。在满足上述满秩条件的点，
该局部内积就是\term{Fisher度量}：它按分布变化而不是坐标差衡量小步长度。对高斯
位置族\(\mathcal{N}(\mu,\sigma^2)\)，得分为\((X-\mu)/\sigma^2\)，故
\[
  F(\mu)=\frac1{\sigma^2}.
\]
局部长度\(|\mathrm{d}\mu|/\sigma\)与KL公式完全一致：噪声越小，相同均值移动在
分布上越显著。
\subsection{自然梯度来自局部分布信赖域}
设要最小化光滑目标\(L(\bm{\theta})\)，梯度为
\(\bm{g}=\nabla L(\bm{\theta})\)。欧氏最速下降取决于当前参数坐标的单位
\scite{math-amari1998natural}若改为限制一步引起的局部KL变化，可解线性化信赖域问题
\begin{equation}
  \min_{\bm{\delta}}
  \bm{g}^{\top}\bm{\delta}
  \quad\text{s.t.}\quad
  \frac12\bm{\delta}^{\top}
  \bm{F}(\bm{\theta})\bm{\delta}
  \leq\varepsilon.
  \label{eq:information-natural-gradient-trust-region}
\end{equation}
当\(\bm{F}\)正定且\(\bm{g}\neq\bm{0}\)时，Lagrange驻点条件
\[
  \bm{g}+\lambda\bm{F}\bm{\delta}=\bm{0}
\]
给出
\begin{equation}
  \bm{\delta}^{\star}
  =
  -
  \sqrt{
    \frac{2\varepsilon}
    {\bm{g}^{\top}\bm{F}^{-1}\bm{g}}
  }\,
  \bm{F}^{-1}\bm{g}.
  \label{eq:information-natural-gradient-step}
\end{equation}
也可以直接验证全局最优性。令
\(\bm{u}=\bm{F}^{1/2}\bm{\delta}\)，则目标为
\((\bm{F}^{-1/2}\bm g)^{\top}\bm u\)，约束为
\(\lVert\bm u\rVert_2\leq\sqrt{2\varepsilon}\)。Cauchy--Schwarz不等式说明
最小值在\(\bm u\)与\(-\bm F^{-1/2}\bm g\)同向且位于边界时达到，代回便得到式
\eqref{eq:information-natural-gradient-step}。若\(\bm g=\bm0\)，线性化目标在
整个可行域上相同，方向不由这个一阶问题唯一确定。
方向\(-\bm{F}^{-1}\bm{g}\)称为\term{自然梯度}（natural gradient）方向。它是
局部二阶近似下的最速下降方向；有限步长下，真实KL与二次近似之间仍有余项。
\subsection{换坐标时为何保持同一切向变化}
令新参数\(\bm{\phi}\)通过光滑可逆映射
\(\bm{\theta}=\bm{\theta}(\bm{\phi})\)给出，Jacobian矩阵为
\(\bm{J}=\partial\bm{\theta}/\partial\bm{\phi}\)。链式法则给出
\[
  \bm{g}_{\bm{\phi}}=\bm{J}^{\top}\bm{g}_{\bm{\theta}},
  \qquad
  \bm{F}_{\bm{\phi}}
  =
  \bm{J}^{\top}\bm{F}_{\bm{\theta}}\bm{J}.
\]
若\(\bm{J}\)可逆且Fisher矩阵正定，则
\[
  \bm{J}\bm{F}_{\bm{\phi}}^{-1}\bm{g}_{\bm{\phi}}
  =
  \bm{F}_{\bm{\theta}}^{-1}\bm{g}_{\bm{\theta}}.
\]
所以两组坐标中的自然梯度对应同一个分布切向方向。这里的不变性是无穷小层面的；具体
离散更新、有限步长和数值近似仍可能随参数化改变。
\subsection{策略几何必须先确定分布对象}
在序贯决策中，“策略的Fisher矩阵”至少可能指三种对象。固定状态\(s\)时，可以对单步
动作分布\(\pi_{\bm{\theta}}(\cdot\mid s)\)取Fisher信息；也可以用指定访问分布
\(d(s)\)平均这些单步矩阵；还可以直接对完整轨迹分布求Fisher信息。三者分别衡量局部
动作变化、访问加权的策略变化和累计轨迹变化。访问分布若也依赖参数，求导时是否将其
视为固定会改变所得对象。
条件高斯模型可以把几种常被混用的矩阵放在同一处比较。设在固定输入分布\(d\)下，
\[
  Y\mid\bm{X}=\bm{x}
  \sim
  \mathcal{N}(\bm{\theta}^{\top}\bm{x},\sigma^2).
\]
单个条件观测的得分、负对数似然Hessian和得分外积分别为
\[
\begin{aligned}
  \bm{s}_{\bm{\theta}}(\bm{x},y)
  &=
  \frac{y-\bm{\theta}^{\top}\bm{x}}{\sigma^2}\bm{x},\\
  -\nabla_{\bm{\theta}}^2\log p_{\bm{\theta}}(y\mid\bm{x})
  &=
  \frac{\bm{x}\bm{x}^{\top}}{\sigma^2},\\
  \bm{s}_{\bm{\theta}}\bm{s}_{\bm{\theta}}^{\top}
  &=
  \frac{(y-\bm{\theta}^{\top}\bm{x})^2}{\sigma^4}
  \bm{x}\bm{x}^{\top}.
\end{aligned}
\]
在模型正确且\(Y\)按该条件分布平均时，残差平方的条件期望为\(\sigma^2\)，后二者
的条件期望相同；再对\(\bm X\sim d\)平均，得到
\(\bm F=\mathbb{E}_{d}[\bm X\bm X^{\top}]/\sigma^2\)。但对一条具体样本，经验
得分外积还乘着随机残差平方，并不等于观测Hessian；模型失配或使用不同输入权重时，
其总体平均也未必等于所定义的Fisher矩阵。
若参数冗余或模型不可识别，\(\bm{F}\)可能奇异。Moore--Penrose伪逆可在可辨识切向
子空间中选取最小范数解，但参数空间中的零空间分量仍不唯一。阻尼
\((\bm{F}+\lambda\bm{I})^{-1}\)保证数值可逆，却加入了坐标相关的欧氏结构；经验
Fisher、分块近似或其他预条件矩阵若不等于真实分布Fisher，也不能无条件继承上述
坐标变换结论。
\section{变分表达与概率学习}
\label{sec:information-variational-learning}
\subsection{证据下界把难积分改写为分布优化}
设联合模型为\(p_{\bm{\theta}}(x,z)\)，观测\(x\)固定，潜变量\(z\)需要积分：
\[
  p_{\bm{\theta}}(x)=\int p_{\bm{\theta}}(x,z)\,\mathrm{d}z.
\]
以下假设\(0<p_{\bm{\theta}}(x)<\infty\)，并取与后验
\(p_{\bm{\theta}}(z\mid x)\)相互绝对连续的概率密度\(q\)；同时假设下列对数比的
正负部分可积。相互绝对连续一方面保证重要性比不会漏掉联合密度的正质量区域，另一方面
保证从\(q\)到后验的KL有定义且有限。将\(q\)乘除进积分，并对凹函数\(\log\)使用
Jensen不等式，
\begin{align}
  \log p_{\bm{\theta}}(x)
  &=\log\mathbb{E}_{q}\!\left[
    \frac{p_{\bm{\theta}}(x,Z)}{q(Z)}\right]\notag\\
  &\geq\mathbb{E}_{q}\!\left[
    \log p_{\bm{\theta}}(x,Z)-\log q(Z)\right].
  \label{eq:information-elbo-jensen}
\end{align}
右侧称为\term{证据下界}（evidence lower bound, ELBO），记为
\(\mathcal{L}(q,\bm{\theta};x)\)。
利用
\(p_{\bm{\theta}}(x,z)=p_{\bm{\theta}}(z\mid x)p_{\bm{\theta}}(x)\)，可精确
分解为
\begin{equation}
  \log p_{\bm{\theta}}(x)
  =\mathcal{L}(q,\bm{\theta};x)+
  D_{\mathrm{KL}}\!\left(
    q(Z)\Vert p_{\bm{\theta}}(Z\mid x)\right).
  \label{eq:information-elbo-decomposition}
\end{equation}
固定\(\bm{\theta}\)时，最大化ELBO等价于在所选变分族中最小化从\(q\)到真实后验的
KL。只有当变分族包含真实后验且优化达到全局最优时，下界才取等号。若同时更新
\(\bm{\theta}\)，证据本身也在改变；近似后验误差、参数拟合误差和优化误差不能合并
成一个“下界差”而不加区分。
若只要求\(q\ll p_{\bm{\theta}}(\cdot\mid x)\)，KL分解仍可在相应有限性条件下使用；
但若后验在\(q\)的零集上还有正质量，前面的重要性采样等号不再成立。反过来，仅让
\(q\)覆盖后验支持也不足以保证反向KL有限。两条推导使用哪一种支撑关系，必须分别核对。

一个线性高斯例子可以核对分解的数值。设
\(Z\sim\mathcal N(0,1)\)、\(X\mid Z\sim\mathcal N(Z,1)\)，观察到\(x=1\)。
此时证据为\(\mathcal N(0,2)\)，后验为\(\mathcal N(1/2,1/2)\)。若取
\(q=\mathcal N(0,1)\)，则
\[
\begin{aligned}
  \log p(x)&=-\frac12\log(4\pi)-\frac14\approx-1.5155,\\
  D_{\mathrm{KL}}\!\left(
    \mathcal N(0,1)\Vert\mathcal N(1/2,1/2)
  \right)
  &=\frac12\left(2+\frac12-1+\log\frac12\right)
  \approx0.4034.
\end{aligned}
\]
因此ELBO约为\(-1.9189\)，与证据对数的差恰为该KL。多变量高斯中，若把\(q\)限制
为分块独立族，同一恒等式仍成立，但不能表达的后验相关性会留在KL近似差距中。
\subsection{Gibbs分布是线性收益与KL代价的最优解}
Donsker--Varadhan表达也可反向写成分布优化。设参考分布为\(P_0\)，并令\(f\)为
有界可测函数。则
\begin{equation}
  \log\mathbb{E}_{P_0}[e^f]
  =\sup_{Q\ll P_0}
  \left\{
    \mathbb{E}_Q[f]-D_{\mathrm{KL}}(Q\Vert P_0)\right\}.
  \label{eq:information-gibbs-variational}
\end{equation}
最优分布满足
\[
  \frac{\mathrm{d}Q^\star}{\mathrm{d}P_0}
  =\frac{e^f}{\mathbb{E}_{P_0}e^f}.
\]
记归一化常数为\(Z_f=\mathbb{E}_{P_0}e^f\)。由于\(f\)有界，\(Q^\star\)下的
\(\mathbb{E}_{Q^\star}f\)和\(D_{\mathrm{KL}}(Q^\star\Vert P_0)\)都有限。对任意
\(Q\ll P_0\)，若\(D_{\mathrm{KL}}(Q\Vert P_0)=+\infty\)，约定目标为
\(-\infty\)；否则
\[
\begin{aligned}
  \mathbb{E}_Q[f]-D_{\mathrm{KL}}(Q\Vert P_0)
  &=
  \log Z_f-D_{\mathrm{KL}}(Q\Vert Q^\star)\\
  &\leq\log Z_f.
\end{aligned}
\]
取\(Q=Q^\star\)达到等号，这便完整给出式~\eqref{eq:information-gibbs-variational}；
因此指数加权不是任意技巧，而是在提高平均收益与偏离参考分布之间作精确权衡的解。
若要允许无界\(f\)，仅有\(0<Z_f<\infty\)还不足以保证\(Q^\star\)取得上确界；例如
\(\mathbb{E}_{P_0}[e^f|f|]<\infty\)是保证上述两项有限的一个充分条件。只有归一化
常数有限而该可积性不成立时，截断倾斜分布仍可逼近上确界，但不能把\(Q^\star\)直接
代入两个无穷项之差。若\(f\)的指数矩无限，归一化常数不存在，右侧优化也可能无有限解。
对有界可测收益\(r\)，加入温度\(\beta>0\)时，
\[
  \sup_Q\left\{
    \mathbb{E}_Q[r]-\beta D_{\mathrm{KL}}(Q\Vert P_0)\right\}
  =
  \beta\log\mathbb{E}_{P_0}[e^{r/\beta}],
\]
且\(Q^\star\propto P_0e^{r/\beta}\)。较小温度更强调高收益区域，但也可能使密度比
集中、采样方差增大；KL正则约束的是相对参考分布的概率变化，不是任务安全约束。
\subsection{指数族与凸对偶}
对指数族
\[
  p_{\bm{\eta}}(x)
  =h(x)\exp\{
    \bm{\eta}^{\top}\bm{T}(x)-A(\bm{\eta})\},
\]
第~\ref{chap:probability-theory}章已得
\(\nabla A(\bm{\eta})=\mathbb{E}_{\bm{\eta}}\bm{T}(X)\)以及
\(\nabla^2A=\operatorname{Cov}_{\bm{\eta}}(\bm{T})\)。对数配分函数的凸共轭为
\begin{equation}
  A^\star(\bm{\tau})
  =\sup_{\bm{\eta}}
  \{\bm{\eta}^{\top}\bm{\tau}-A(\bm{\eta})\}.
  \label{eq:information-log-partition-conjugate}
\end{equation}
若\(\bm{\tau}\)位于可实现均值参数集合的相对内部，且存在
\(\bm{\eta}\)满足\(\nabla A(\bm{\eta})=\bm{\tau}\)，则最优性条件给出
\[
  A^\star(\bm{\tau})=\bm{\eta}^{\top}\bm{\tau}-A(\bm{\eta}).
\]
相对于基准密度\(h\)的负熵正是这一共轭量。于是自然参数\(\bm{\eta}\)与均值参数
\(\bm{\tau}\)形成对偶坐标，Fisher矩阵
\(\nabla^2A(\bm{\eta})\)把两者的局部变化联系起来。
若指数族不是最小表示，充分统计量之间存在线性冗余，\(A\)不会严格凸，均值参数与自然
参数也不一一对应。边界均值可能只能由参数趋于无穷的极限实现。最大熵与凸对偶结论必须
在可行域及其相对内部中解释。
例如在有限集合\(\{0,1,2\}\)上要求\(\mathbb{E}X=1/2\)，相对于均匀计数测度的
最大熵分布满足\(p_j\propto t^j\)。矩约束给出
\[
  \frac{t+2t^2}{1+t+t^2}=\frac12,
  \qquad
  t=\frac{\sqrt{13}-1}{6}\approx0.4343,
\]
从而\((p_0,p_1,p_2)\approx(0.6162,0.2676,0.1162)\)。这里优化变量是满足矩约束
的分布，\(t=e^\lambda\)只是实现约束的乘子参数；若改为拟合一个既定数据分布，
优化对象和保证都已改变。
\subsection{学习目标中的变分量不等于概率保证}
知识蒸馏常让学生分布\(Q_{\bm{\phi}}(\cdot\mid x)\)逼近教师分布
\(P_{\mathrm{T}}(\cdot\mid x)\)。最小化
\(D_{\mathrm{KL}}(P_{\mathrm{T}}\Vert Q_{\bm{\phi}})\)等价于教师软标签下的交叉熵，
强调覆盖教师赋正质量的类别；反向KL会采用不同的平均分布与支持要求。无论哪一方向，
学生得到的是对教师分布的逼近，不自动比教师更接近真实标签分布。
对比学习常把一个配对样本与若干负样本组成分类任务。适当采样条件下，分类对数比可以
形成互信息的变分下界；但界的对象取决于正负样本的联合与边缘分布，有限负样本数还会
限制界的紧度。优化某个对比目标不能在未核对采样机制时直接解释为“最大化全部语义
信息”。
比较隐私机制输出时，KL或R\'enyi散度可以作为变分控制中的代价；正式隐私保证还要求
对所有相邻输入统一成立，并按机制组合规则换算。类似地，ELBO较高只说明所选联合模型
和变分族下的下界较高，不等于后验已经校准，也不等于生成样本满足任务评价。变分表达
把难问题转成优化问题，但不会替代模型假设、统计证据和任务保证。
\section{本章小结}
\label{sec:information-theory-statistical-geometry-summary}
熵从平均描述长度刻画单个分布的不确定性，条件熵和互信息说明已有信息与新增观测分别
减少多少不确定性。对数评分把预测分布放入真实分布下评价，交叉熵由真实熵与KL失配
组成；绝对连续性决定这种失配是否有限，KL的方向决定重点覆盖哪些概率区域。
分布差异没有唯一答案。总变差控制事件和有界测试函数，Hellinger提供对称的平方根密度
距离，MMD只考察所选核函数类，R\'enyi散度强调似然比高阶矩。Wasserstein距离进一步
把底层空间的搬运成本纳入比较，因而能连续反映支持位置的移动。选择差异时必须说明要
控制的事件、损失、矩或空间成本，不能把一个距离的小值无条件替换成另一个距离的小值。
换测度变分式把指数矩与KL代价连接起来，可用于信息稳定性、PAC--Bayes和轨迹比较。
Le Cam与Fano从相反方向说明，如果目标彼此分离而观测分布仍接近，任何估计器都必须
承担误差；自适应访问下应比较算法诱导的完整轨迹，不能假装获得了固定设计的独立样本。
在得分满秩的正则模型中，KL的局部二阶项给出Fisher度量，自然梯度据此限制一步造成的
分布变化。奇异模型、阻尼与近似预条件会改变这项几何保证。ELBO、Gibbs变分式和指数族凸
对偶则说明，许多概率学习目标都可看成在期望收益、熵与相对熵之间优化；优化替代分布
并不自动建立真实世界的概率保证。信息不足能够限制估计，但观测分布能否回答干预问题，
还需要第~\ref{chap:causal-inference}章进一步引入生成机制与识别假设。
